\documentclass[12pt]{article}
% Préambule commun à tous les documents extraits de « La Revue CAPES »
\usepackage[T1]{fontenc}
\usepackage[utf8]{inputenc}
\usepackage{lmodern}
\usepackage{amsmath,amssymb,amsthm,amsfonts,mathrsfs}
\usepackage{setspace}
\usepackage{listings}
\usepackage{pgfplots}
\pgfplotsset{compat=1.18}
\usepackage{longtable}
\usepackage{adjustbox}
\usepackage{lettrine}
\usepackage{tkz-tab}
\usepackage{changepage}
\usepackage{pdflscape}
\usepackage{graphicx}
\usepackage{tikz}
\usepackage{xcolor}
\usepackage[a4paper,top=2cm,bottom=2.5cm,left=2.5cm,right=2cm]{geometry}
\usepackage{fancyhdr}
\usepackage{draftwatermark}
\SetWatermarkText{La RevueCapes}
\SetWatermarkLightness{0.9}
\SetWatermarkScale{2}
\usepackage{hyperref}
\hypersetup{colorlinks=true,linkcolor=blue!50!black,urlcolor=blue!50!black}

\definecolor{revue}{RGB}{20,60,120}

\pagestyle{fancy}
\fancyhf{}
\renewcommand{\headrulewidth}{0.4pt}
\setlength{\headheight}{15pt}

% \entete{Matière}{Titre}{Type de document}
\newcommand{\entete}[3]{%
  \fancyhead[L]{\small\textsc{La Revue CAPES} -- #1}%
  \fancyhead[R]{\small #2}%
  \fancyfoot[C]{\thepage}%
  \thispagestyle{plain}%
  \begin{center}
    {\color{revue}\rule{\textwidth}{1.2pt}}\\[4pt]
    {\small\textsc{Concours directs CAP-CEG / CAPES -- Burkina Faso}}\\[6pt]
    {\LARGE\bfseries\color{revue} #2}\\[6pt]
    {\large\itshape #1 -- #3}\\[2pt]
    {\color{revue}\rule{\textwidth}{1.2pt}}
  \end{center}
  \vspace{0.5cm}%
}

% Encadré pour les remarques ajoutées dans les corrigés
\newcommand{\remarque}[1]{\par\medskip\noindent\fcolorbox{revue}{revue!6}{\parbox{\dimexpr\linewidth-2\fboxsep-2\fboxrule}{\textbf{\color{revue}Remarque.} #1}}\par\medskip}


\begin{document}
\entete{Culture générale}{Corrigé du sujet type de dissertation n°1}{Correction}
\begin{spacing}{1.5}

\textbf{ Introduction}

\textbf{ Accroche :} La poésie est souvent perçue comme un genre littéraire qui plonge au cœur des émotions humaines, explorant des thèmes profonds comme la douleur, la mélancolie, ou les tourments de l'âme.

\textbf{ Présentation de la citation :} \emph{Roland Barthes}, dans Le degré zéro de l'écriture, affirme que l'univers poétique est habité par des tourments, suggérant que les poètes sont marqués par une profonde gravité, incapables de sourire.

\textbf{ Problématique :} Cette vision des poètes comme figures tourmentées est-elle représentative de l'ensemble de l'univers poétique ? Existe-t-il des contre-exemples ou des nuances à apporter à cette idée ?

\textbf{ Annonce du plan :} Nous examinerons tout d'abord les œuvres qui confirment cette opinion, avant d'explorer les exceptions et de conclure sur la diversité de l'univers poétique.

\quad
I. L'univers poétique comme miroir des tourments humains
Barthes semble mettre en avant la capacité de la poésie à exprimer les souffrances et les tourments, qui sont des thèmes récurrents dans l'histoire littéraire.

1. Les poètes maudits : une figure emblématique de la souffrance poétique

Exemple : Charles Baudelaire (Les Fleurs du mal)
Baudelaire explore des thèmes comme le spleen, l'ennui, et la déchéance. Ses poèmes, tels que L'Albatros ou Spleen, mettent en scène un univers sombre où le poète est en proie à une souffrance existentielle.
Paul Verlaine et les tourments de l'amour (Romances sans paroles, Sagesse).
Son écriture reflète des conflits intérieurs liés à l'amour, la culpabilité, et la recherche spirituelle.

2. La poésie comme exutoire des douleurs personnelles

Exemple : Arthur Rimbaud (Une Saison en enfer).
Rimbaud exprime sa révolte contre le monde et ses propres blessures intérieures.
Anna Akhmatova : Dans sa poésie, la douleur liée à la guerre et à la répression soviétique imprègne ses œuvres, notamment dans Requiem.

\quad
II. Une vision partielle : la poésie peut aussi célébrer la joie et l'émerveillement
Bien que Barthes souligne les tourments des poètes, il existe des œuvres poétiques qui explorent des thèmes légers, joyeux ou empreints d'émerveillement.

1. La célébration de la beauté et de la nature

Exemple : Alphonse de Lamartine (Le Lac).
Bien que teinté de mélancolie, ce poème célèbre également la beauté de la nature et l'immortalité des souvenirs.
Victor Hugo (Les Contemplations).
Certains poèmes de Hugo, comme Demain, dès l'aube, mêlent douleur et espoir, montrant que l'écriture poétique n'est pas uniquement marquée par la souffrance.

2. L'humour et l'ironie dans la poésie

Exemple : Jacques Prévert (Paroles).
Prévert écrit des poèmes pleins d'ironie et de légèreté, jouant avec les mots et les images pour offrir un regard critique mais amusé sur la société.
Francis Ponge : Dans Le Parti pris des choses, Ponge célèbre des objets simples de la vie quotidienne avec une poésie décalée et ludique.

\quad
III. Vers une vision nuancée de la condition du poète
La citation de Barthes peut être interprétée comme une généralisation. Cependant, l'univers poétique est divers et reflète une gamme complète d'émotions humaines.

1. La poésie comme catharsis : une nécessité d'exprimer les tourments
Même lorsque la poésie aborde des thèmes sombres, elle permet au poète de se libérer de ses souffrances. Cela peut être vu comme une forme de sourire intérieur, une réconciliation avec soi-même.

2. La poésie moderne : un espace d'exploration ludique et expérimental
Dans la poésie contemporaine, de nombreux poètes, comme Raymond Queneau (Cent Mille Milliards de Poèmes), utilisent l'écriture poétique pour expérimenter et jouer avec la langue, rompant avec l'image traditionnelle du poète tourmenté.

\quad
\quad
\textbf{ Conclusion}
\textbf{Résumé : }L'affirmation de Barthes souligne à juste titre la prévalence des thèmes sombres et tourmentés dans l'univers poétique, comme en témoignent des œuvres de Baudelaire, Verlaine ou Rimbaud. Cependant, cette vision n'englobe pas toute la richesse de la poésie, qui peut aussi célébrer la joie, l'émerveillement, et l'humour.

\textbf{Ouverture :} La poésie, par sa diversité, reflète toute la complexité de la condition humaine. Peut-être que l'univers poétique n'est pas tant un lieu de souffrance qu'un espace où les émotions humaines, qu'elles soient sombres ou lumineuses, trouvent une expression universelle.




\quad

\end{spacing}
\end{document}
